% ============================================================
% sec/02related_work.tex
% Core section of Stage 1. Three mandatory axes + comparative
% table + literature baseline + explicit gap identification.
% Minimum 21 peer-reviewed references (7 per member).
% ============================================================
\section{Related Work}
\label{sec:related}

\todo{One short paragraph stating the review methodology: databases
queried, search strings, inclusion and exclusion criteria, time window,
and how many works were retained out of how many screened.}

\subsection{Axis A: The Problem Domain}
\label{subsec:axisa}
\emph{How has this problem been addressed in the literature? Which
approaches dominate? What is considered solved and what remains open?}

\todo{3--5 paragraphs synthesizing at least 7 peer-reviewed works.
Group by approach family, not one paragraph per paper.}

\todo{Closing paragraph: which approaches dominate, what is settled,
what is still open.}

\subsection{Axis B: Techniques of the Selected Track}
\label{subsec:axisb}
\emph{What is the current state of the methods we intend to use? What
are their strengths, limits, and the assumptions under which they work?}

\todo{3--5 paragraphs synthesizing at least 7 peer-reviewed works on the
technique family of the chosen track.}

\todo{Closing paragraph: the technical justification of the track,
supported by literature and not by preference. State explicitly which
alternative track or technique family was discarded and why.}

\subsection{Axis C: Dataset / Environment}
\label{subsec:axisc}
\emph{Which works have used it? Which metrics do they report and under
which protocol? What is the best published result? Which known problems
does it have?}

\todo{3--5 paragraphs synthesizing at least 7 peer-reviewed works that
use the selected dataset or environment.}

\todo{Documented known problems: bias, leakage, size, representativeness.
For an RL environment: version differences, reward scaling, and
non-comparable evaluation protocols.}

\subsection{Comparative Analysis}
\label{subsec:comparative}
\safeinput{tab/comparative_table}

\todo{2--3 paragraphs reading the table: which trends are visible,
which comparisons are unfair and why, which limitation recurs across
works.}

\subsection{Literature Baseline}
\label{subsec:litbaseline}
\todo{State the best published result on the selected dataset or
environment under the selected metric, with its reference. If a direct
comparison is impossible, state a defensible range instead.}

\begin{quote}
\textbf{Target:} \todo{The number this work aims to match or exceed,
with its metric and its source.}
\end{quote}

\noindent\textbf{Protocol differences that block a fair comparison.}
\begin{itemize}
  \item \todo{Different train/test split}
  \item \todo{Different metric or metric variant}
  \item \todo{Different data subset or environment version}
  \item \todo{Different amount of training budget or compute}
\end{itemize}

\subsection{Identified Gaps}
\label{subsec:gaps}
\textbf{Gaps found in the literature.}
\begin{enumerate}
  \item \todo{Gap 1}
  \item \todo{Gap 2}
  \item \todo{Gap 3}
\end{enumerate}

\noindent\textbf{Gaps this work addresses.} \todo{Which of the above,
and how.}

\noindent\textbf{Gaps out of scope.} \todo{Which ones remain unaddressed
and why they exceed the scope of this project.}